% 08-gcli.tex — GCLI、预测与显著性模型
% Source: chapters/08-gcli-prediction-and-significance-model.md (expanded)

\chapter{GCLI、预测与显著性模型}\label{ch:gcli}

\section{本章目标}

\begin{itemize}
  \item 理解 GCLI 存在的原因：它解决的是``每组系数大概有多大''的摘要问题
  \item 掌握 GCLI 的计算：OR $\to$ BSR $\to$ GCLI
  \item 明确区分三个容易混淆的概念：coefficient、GCLI、GTLI
  \item 理解 GCLI 编码的三层机制：预测残差 $\to$ 显著性标志 $\to$ unary 编码
  \item 掌握 GCLI 方法的选择机制（signaling bits）
\end{itemize}

\section{GCLI 是什么——先建立直觉}

\subsection{它想解决什么问题}

JPEG XS 不使用算术编码。它的``熵编码''部分非常轻量，核心就是 \textbf{GCLI 编码}。

要理解 GCLI 为什么存在，先想一个问题：编码器已经拿到了一组系数（DWT 变换、sign-magnitude 表示），接下来要按 bitplane 逐层编码。但是——

\begin{tcolorbox}[colback=yellow!5,colframe=yellow!60!black,title=核心问题]
编码器怎么知道该从哪一层 bitplane 开始编码？
\end{tcolorbox}

如果一组 4 个系数是 `[3, 7, 2, 5]`，它们中最大的是 7（= \texttt{0b111}），最高有效位在 bitplane 2。那编码器只需要编码 bitplane 2、1、0 就够了——bitplane 3 以上全是零，编码纯属浪费。

\textbf{GCLI 就是回答这个问题的摘要信息}。它不是编码``这 4 个系数分别是多少''，而是编码``这 4 个系数最高用到了第几位''。

你可以先把 GCLI 理解成这一组系数的``位深标签''：它告诉后续编码模块，这一组最多需要关心哪些 bitplane。

\subsection{一个极小例子}

假设一组 4 个系数的幅度值为 $[3, 7, 2, 5]$。

\textbf{第 1 步：按位 OR 聚合}

把 4 个数的二进制表示做按位 OR：

\begin{lstlisting}[numbers=none]
3 = 0b0011
7 = 0b0111
2 = 0b0010
5 = 0b0101
OR  -------
    0b0111 = 7
\end{lstlisting}

OR 的作用是：只要任意一个系数在某个 bitplane 上是 1，结果在该位上就是 1。所以 OR 结果保留了这组系数中``用到的所有 bitplane''的信息。

\textbf{第 2 步：找最高非零位（BSR）}

$7 = \texttt{0b111}$，最高非零位是 bit 2。$\mathrm{BSR}(7) = 2$。

\textbf{第 3 步：GCLI = BSR + 1}

$\mathrm{GCLI} = 2 + 1 = 3$。

这表示：这组系数最多需要 3 个 bit（bitplane 2、1、0）来表示。

\textbf{第 4 步：与 bitplane 编码的关系}

知道了 GCLI = 3，编码器就知道：
\begin{itemize}
  \item bitplane 3 及以上：全是零，不需要编码
  \item bitplane 2、1、0：需要编码
\end{itemize}

如果后续量化参数 GTLI = 1（第 10 章会详细讲），那么实际只编码 bitplane 2 和 1（即 $\mathrm{GCLI}-1$ 到 $\mathrm{GTLI}$），bitplane 0 被截断丢弃。

\subsection{三个容易混淆的概念}

在继续之前，请先明确区分以下三个概念。初学者最容易在这里犯晕：

\begin{table}[H]
\centering
\caption{三个容易混淆的概念}
\begin{tabular}{@{}lp{4.5cm}lp{3cm}@{}}
\toprule
概念 & 它描述什么 & 谁产生的 & 类比 \\
\midrule
\textbf{coefficient}（系数） & 每个系数的实际值 & DWT 变换输出 & 货物本身 \\
\textbf{GCLI} & 一组系数的最高位深（摘要） & 对系数做 OR + BSR & 货物的``体积等级'' \\
\textbf{GTLI} & 量化时截断到哪一位 & 码率控制决定 & ``只保留到哪一级''的截断指令 \\
\bottomrule
\end{tabular}
\end{table}

关系：
\begin{itemize}
  \item \textbf{GCLI 是系数的属性}——从系数算出来的，描述``这组系数有多大''
  \item \textbf{GTLI 是控制参数}——由码率控制决定的，描述``预算不够了，只能保留到哪一位''
  \item 后续 bitplane 编码的范围是 $\mathrm{GCLI}-1$ 到 $\mathrm{GTLI}$——GCLI 决定上界，GTLI 决定下界
\end{itemize}

\begin{tcolorbox}[colback=yellow!5,colframe=yellow!60!black,title=关键区分]
GCLI 告诉你``这组系数理论上最高用到哪一层''，GTLI 告诉你``实际编码时保留到哪一层''。GCLI 来自数据，GTLI 来自预算。
\end{tcolorbox}

\section{数学定义}

\subsection{GCLI 计算}

给定一个系数组的 4 个 sign-magnitude 值，GCLI 的计算过程：
\begin{equation}
\mathit{GCLI} = \begin{cases} 0 & \text{if } \mathit{or\_all} = 0 \\ \lfloor \log_2(\mathit{or\_all}) \rfloor + 1 & \text{otherwise} \end{cases}
\end{equation}

其中 $\mathit{or\_all} = v_0 \mid v_1 \mid v_2 \mid v_3$（取幅度部分，不含符号位）。

等价于：GCLI 是该组中最大系数的二进制位数。

在代码中（\texttt{precinct.c:87-89}）：
\begin{lstlisting}[language=C,numbers=none]
#define BSR(x) (31 - __builtin_clz(x))  // bit scan reverse
#define GCLI(x) (((x) == 0) ? 0 : (BSR((x)) + 1))
\end{lstlisting}

\subsubsection{源码摘录：\texttt{compute\_gcli\_buf()}}

\begin{lstlisting}[language=C,numbers=none]
int compute_gcli_buf(const sig_mag_data_t* in, int len,
                     gcli_data_t* out, int max_out_len,
                     int* out_len, int group_size)
{
    sig_mag_data_t or_all;
    *out_len = (len + group_size - 1) / group_size;

    // (1) 完整组（每组 group_size=4 个系数）
    for (i = 0; i < (len / group_size) * group_size; i += group_size)
    {
        for (or_all = 0, j = 0; j < group_size; j++)
            or_all |= *in++;                                  // OR 聚合
        *out++ = GCLI(or_all & (~SIGN_BIT_MASK));             // BSR+1
    }

    // (2) 尾部不足一组
    if (len % group_size)
    {
        for (or_all = 0, j = 0; j < len % group_size; j++)
            or_all |= *in++;
        *out++ = GCLI(or_all & (~SIGN_BIT_MASK));
    }
    return 0;
}
\end{lstlisting}

\begin{engineeringnote}
\textbf{逐段解释}：
\begin{itemize}
  \item \textbf{(1) 完整组}：每 4 个系数为一组，对幅度部分（\texttt{\& \~{}SIGN\_BIT\_MASK} 屏蔽符号位）做按位 OR，然后取 $\mathrm{BSR} + 1$。OR 聚合确保只要任一系数在某个 bitplane 上非零，该 bitplane 就被保留。
  \item \textbf{(2) 尾部处理}：当系数总数不是 4 的倍数时，最后一组不满 4 个。处理方式相同，只是 OR 的操作数更少。
\end{itemize}
\end{engineeringnote}

\subsubsection{源码摘录：\texttt{update\_gclis()}}

\begin{lstlisting}[language=C,numbers=none]
void update_gclis(precinct_t* prec)
{
    int band, ypos;
    for (band = 0; band < bands_count_of(prec); band++)
    {
        for (ypos = 0; ypos < precinct_in_band_height_of(prec, band); ypos++)
        {
            const sig_mag_data_t* in = precinct_line_of(prec, band, ypos);
            gcli_data_t* dst = precinct_gcli_of(prec, band, ypos);
            int reclen, width = (int)precinct_width_of(prec, band);
            int gcli_count = (int)precinct_gcli_width_of(prec, band);
            compute_gcli_buf(in, width, dst, gcli_count, &reclen,
                             prec->group_size);
        }
    }
}
\end{lstlisting}

遍历所有 band 的所有行，逐行调用 \texttt{compute\_gcli\_buf()}。

\subsection{完整概念链：从系数到编码}

从原始系数到最终编码输出，数据经过以下变换。请按顺序理解每一层：

\begin{table}[H]
\centering
\small
\caption{完整概念链}
\begin{tabular}{@{}cL{2.3cm}L{2.8cm}L{5.2cm}@{}}
\toprule
步骤 & 概念 & 含义 & 示例 \\
\midrule
1 & \textbf{magnitude} & sign-magnitude 去掉符号位的幅度值 & 系数 \texttt{0x80000005} $\to$ 幅度 = 5 \\
2 & \textbf{OR} & 同组 4 个系数幅度的按位或 & $[3, 5, 0, 2] \to 3\,|\,5\,|\,0\,|\,2 = 7$ \\
3 & \textbf{BSR} & OR 结果的最高非零位位置 & $\mathrm{BSR}(7) = 2$（$7 = \texttt{0b111}$） \\
4 & \textbf{GCLI} & BSR + 1，即该组位深 & $\mathrm{GCLI} = 3$ \\
5 & \textbf{predictor} & 上方 precinct 同位置的 GCLI 值 & pred = 4（上一行同 band 同组） \\
6 & \textbf{residual} & GCLI $-$ predictor（预测残差） & $3 - 4 = -1$ \\
7 & \textbf{sigflag} & 该组是否需要编码（非零残差/非零系数） & residual $\neq 0$ $\to$ sigflag = 1 \\
\bottomrule
\end{tabular}
\end{table}

最终编码流程：magnitude $\to$ OR $\to$ BSR+1 = GCLI $\to$ 减去 predictor 得 residual $\to$ 若 sigflag = 1，对 residual 做 unary 编码。

\section{预测与残差——为什么不是直接编码 GCLI}

\subsection{预测的直觉}

GCLI 值本身可以直接编码（每个值 4 bit），但 JPEG XS 做了一个优化：\textbf{用上方 precinct 同位置的 GCLI 作为预测值，只编码差值}。

\begin{engineeringnote}
相邻 precinct（上下方向）的同一 band 同一组位置，其系数统计特性往往相似——上方用了 4 个 bitplane 的位置，当前 precinct 大概率也差不多。只编码一个小的差值（residual）比编码完整值更省 bit。
\end{engineeringnote}

\subsection{预测类型}

\begin{table}[H]
\centering
\caption{预测类型}
\begin{tabular}{@{}cll@{}}
\toprule
预测类型 & 符号 & 含义 \\
\midrule
无预测 & PRED\_NONE & 直接编码 GCLI 值 \\
垂直预测 & PRED\_VER & 用上方 precinct 同组的 GCLI 作为预测 \\
\bottomrule
\end{tabular}
\end{table}

垂直预测的计算（\texttt{pred.c}）：
\begin{equation}
\mathit{residual}[i] = \mathit{gcli\_current}[i] - \mathit{gcli\_top}[i]
\end{equation}

编码器传输 \texttt{residual}，解码器加上 \texttt{gcli\_top} 恢复原始 GCLI。

\subsection{残差为零时的优化}

如果预测残差恰好为 0（上方 GCLI 和当前 GCLI 相同），那就不需要编码任何东西。这就是\textbf{显著性标志}（sigflag）发挥作用的地方——它先告诉解码器``这一组的残差是不是零''，零的直接跳过。

\section{显著性标志（Significance Flags）}

\subsection{为什么需要 sigflag}

编码器逐行处理 GCLI 值（对预测编码而言是 GCLI 残差）。一行中很多值可能为零——预测准确时残差为零，或者该位置本身就没有有效系数。sigflag 的作用就是在正式编码 GCLI 之前，用一组粗粒度的标志位标记哪些位置``需要继续看''，零值位置直接跳过。

\subsection{sigflag 的粒度：$S_s$ 参数}

sigflag 不是为每个 GCLI 值（或每个 GCLI 组）单独生成一个标志位。它把 GCLI 值按 \textbf{$S_s$ 个为一组}（$S_s$ 默认为 8），每组只产生 1 个 sigflag bit。如果这一组内任何一个 GCLI 值非零，标志位就是 1；全零则标志位为 0。

源码中（\texttt{sig\_flags.c}），sigflag 的层级数量由下式决定：
\begin{equation}
\mathit{lvls\_size} = \lceil \mathit{buf\_len} \,/\, S_s \rceil
\end{equation}

即一行有 $N$ 个 GCLI 值时，sigflag 只需要 $\lceil N / S_s \rceil$ bit。这与 GCLI 计算的组大小 $N_g = 4$ 是两个独立的参数——$N_g$ 控制 GCLI 的计算粒度（4 个系数共享一个 GCLI），$S_s$ 控制 sigflag 的编码粒度（$S_s$ 个 GCLI 值共享一个 sigflag）。

\subsection{两种 Run 模式}

sigflag 的判断依据取决于 run 模式：

\begin{table}[H]
\centering
\caption{两种 Run 模式}
\begin{tabular}{@{}clcp{5cm}@{}}
\toprule
模式 & 符号 & $R_m$ 值 & sigflag 判断依据 \\
\midrule
ZRF & RUN\_SIGFLAGS\_ZRF & 0 & 该 $S_s$ 个 GCLI 残差中是否有非零值 \\
ZRCSF & RUN\_SIGFLAGS\_ZRCSF & 1 & 该 $S_s$ 个 GCLI 原始值中是否有非零值 \\
\bottomrule
\end{tabular}
\end{table}

\begin{engineeringnote}
ZRF 和 ZRCSF 的区别仅在于\textbf{sigflag 的输入数据不同}（见 \texttt{gcli\_budget\_compute\_generic()}，\texttt{gcli\_budget.c:135-136}）：
\begin{itemize}
  \item \textbf{ZRF}：sigflag 的输入是预测残差（residuals），判断``这 $S_s$ 个残差里有没有非零的''
  \item \textbf{ZRCSF}：sigflag 的输入是 GCLI 原始值（values，即无预测的 GCLI），判断``这 $S_s$ 个 GCLI 值里有没有非零的''
\end{itemize}
sigflag 过滤后，非零位置对应的 GCLI 残差（ZRF）或 GCLI 原始值（ZRCSF）才进入后续编码。
\end{engineeringnote}

\subsection{sigflag 编码的完整过程}

用一个具体例子说明。假设一个 band 行有 8 个 GCLI 残差，$S_s = 4$，使用 ZRF 模式。

\begin{center}
\small
\begin{tabular}{@{}ccccccccc@{}}
\toprule
位置 & 0 & 1 & 2 & 3 & 4 & 5 & 6 & 7 \\
\midrule
残差 & $0$ & $+2$ & $0$ & $-1$ & $0$ & $0$ & $+3$ & $0$ \\
\bottomrule
\end{tabular}
\end{center}

\textbf{Step 1：按 $S_s = 4$ 分组并生成 sigflag}

\begin{lstlisting}[numbers=none]
S_s = 4, lvls_size = ceil(8/4) = 2 个 sigflag bit

Chunk 0 (位置 0-3): 残差 [0, +2, 0, -1] → 有非零 → sigflag[0] = 1
Chunk 1 (位置 4-7): 残差 [0, 0, +3, 0]  → 有非零 → sigflag[1] = 1
\end{lstlisting}

sigflag 段输出：\texttt{1 1}（2 bit）

\textbf{Step 2：写入码流}

sigflag 以每个 1 bit 写入（\texttt{sig\_flags\_write()} 中逐位 \texttt{bitpacker\_write(bitstream, val, 1)}）。写入的是标志位的逻辑反（源码做了 \texttt{val = !val}，这是编码约定，解码端对应反转），但不影响 bit 数——就是 $\mathit{lvls\_size}$ bit。

\textbf{Step 3：过滤 GCLI 编码}

解码器读到 sigflag = \texttt{[1, 1]}（两个 chunk 都有非零值），所以两个 chunk 的所有残差都进入后续 GCLI 编码。本例中所有 8 个残差都被编码——因为每个 chunk 至少有一个非零值，两个 sigflag bit 都是 1。

反之，如果某个 chunk 全为零（sigflag = 0），该 chunk 的所有残差被跳过，不进入 GCLI unary 编码。这就是 sigflag 节省 bit 的场景。

\textbf{sigflag 预算}：\texttt{sig\_flags\_budget()} 返回 $\mathit{lvls\_size} = \lceil N / S_s \rceil$。上例中 $N = 8, S_s = 4$，预算为 2 bit。如果 $S_s = 8$ 且 $N = 8$，则 $\mathit{lvls\_size} = \lceil 8/8 \rceil = 1$，仅需 1 bit。

\subsection{为什么不是每组 1 bit}

初学者容易把 $S_s$ 和 $N_g$ 混淆。$N_g = 4$ 是 GCLI 计算的组大小——4 个 sign-magnitude 系数通过 OR + BSR 产生 1 个 GCLI 值。$S_s = 8$ 是 sigflag 的组大小——8 个 GCLI 值（对应 $8 \times 4 = 32$ 个原始系数）共享 1 个 sigflag bit。

这个设计的工程理由：sigflag 越细粒度，跳过的机会越多，但 sigflag 本身的 bit 开销也越大。$S_s = 8$ 是在大多数场景下的工程权衡——一个 sigflag bit 覆盖 32 个系数的编码决策。

\section{编码方法}

\subsection{Unary 编码}

Unary 编码是最简单的变长编码：用一串 0 后跟一个 1 来表示非负整数 $n$：

\begin{equation}
n = 0 \to \texttt{1}, \quad n = 1 \to \texttt{01}, \quad n = 2 \to \texttt{001}, \quad \ldots
\end{equation}

对于有符号的预测残差，使用 4-clipped signed alphabet（将正负值交替映射到编码值）：

\begin{table}[H]
\centering
\caption{4-clipped signed alphabet}
\begin{tabular}{@{}cll@{}}
\toprule
预测残差 & 编码（0 的个数 + 终止 1） & 位数 \\
\midrule
0 & \texttt{1} & 1 \\
$-1$ & \texttt{01} & 2 \\
$+1$ & \texttt{001} & 3 \\
$-2$ & \texttt{0001} & 4 \\
$+2$ & \texttt{00001} & 5 \\
$\pm k$ ($3 \leq k \leq 12$) & $\texttt{0}^{k+3}\texttt{1}$ & $k + 4$ \\
$|\Delta| \geq 13$ & $\texttt{0}^{15}\texttt{1}$ & 16（最大） \\
\bottomrule
\end{tabular}
\end{table}

\begin{engineeringnote}
残差为 0 是最常见的情况（预测准确），所以编码最短（1 bit）。残差绝对值越大越少见，编码越长。这是一个自然的概率匹配。
\end{engineeringnote}

\subsubsection{编码示例}

假设经过 sigflag 过滤后，需要编码的残差序列为 $[0, +1, -2, 0, +3]$。使用 4-clipped signed alphabet：

\begin{lstlisting}[numbers=none]
残差值 → 编码       → 位数
  0    → 1          → 1 bit
 +1    → 001        → 3 bit
 -2    → 0001       → 4 bit
  0    → 1          → 1 bit
 +3    → 0000001    → 7 bit
\end{lstlisting}

完整 bit 流（串联）：\texttt{1 001 0001 1 0000001}，共 16 bit。

解码规则：读取连续的 0，遇到 1 停止。数 0 的个数 $n$：
\begin{itemize}
  \item $n = 0$：值 $= 0$
  \item $n$ 为奇数：值 $= -(n+1)/2$
  \item $n$ 为偶数（$>0$）：值 $= +n/2$
\end{itemize}

解码过程：

\begin{lstlisting}[numbers=none]
读到 1       → 0 个零 → 值 = 0          ✓
读到 001     → 2 个零 → 值 = +2/2 = +1  ✓
读到 0001    → 3 个零 → 值 = -(4/2) = -2 ✓
读到 1       → 0 个零 → 值 = 0          ✓
读到 0000001 → 6 个零 → 值 = +6/2 = +3  ✓
\end{lstlisting}

\subsubsection{源码摘录：\texttt{unary\_encode()}}

\begin{lstlisting}[language=C,numbers=none]
int unary_encode(bit_packer_t* bitstream, gcli_pred_t* gcli_pred_buf,
                 int len, int no_sign, unary_alphabet_t alph)
{
    int n = 1, bit_len = 0;
    for (i = 0; (i < len) && (n > 0); i++)
    {
        if (!no_sign)
            n = bitpacker_write_unary_signed(bitstream, gcli_pred_buf[i], alph);
        else
            n = bitpacker_write_unary_unsigned(bitstream, gcli_pred_buf[i]);
        bit_len += n;
    }
    return bit_len;
}
\end{lstlisting}

有预测时（\texttt{no\_sign = 0}）使用带符号 unary 编码残差；无预测时（\texttt{no\_sign = 1}）使用无符号 unary 编码 GCLI 值本身。

\subsection{Raw 4-bit 编码}

Raw 编码是最简单的方式：每个 GCLI 值直接用固定的 4 bit 写入（GCLI 范围 0--15，恰好需要 4 bit）。没有预测，没有 sigflag，没有 unary 变长编码。

它并不依靠启发式判断``分布是否均匀''来决定是否使用。实际机制是 \texttt{detect\_gcli\_raw\_fallback()} 在 subpacket 层面对比两种方案的总 bit 数：如果 raw（4 bit $\times$ GCLI 个数）$\leq$ sigflag + unary 编码的总和，就用 raw 替代。这是一种精确的预算比较，详见第 \ref{ch:rate-control} 章。

\subsection{Bounded 编码}

Bounded 编码是 unary 编码的优化变体。它的核心思想：\textbf{利用 GTLI 信息排除不可能的残差值，缩小编码范围，从而减少平均编码长度}。

回忆一下：GCLI 的取值范围是 $[0, 15]$。但有了 GTLI 和 predictor 以后，很多残差值实际上不可能出现：
\begin{itemize}
  \item GCLI 不能低于 GTLI（低于 GTLI 的组在量化阶段就被清零了，不会走到 GCLI 编码），所以 $\mathrm{GCLI} \geq \mathrm{GTLI}$
  \item GCLI 不能高于 $\mathrm{MAX\_GCLI} = 15$
  \item 结合以上两点，残差 $\Delta = \mathrm{GCLI} - \mathit{predictor}$ 的真实范围取决于 predictor 和 GTLI 的相对大小
\end{itemize}

源码中 \texttt{bounded\_code\_get\_min\_max()} 计算了两个边界：
\begin{align}
\mathit{min\_allowed} &= -\max(\mathit{predictor} - \mathit{GTLI},\; 0) \label{eq:bounded-min}\\
\mathit{max\_allowed} &= \max(15 - \max(\mathit{predictor},\, \mathit{GTLI}),\; 0) \label{eq:bounded-max}
\end{align}

理解这两个公式的关键：
\begin{itemize}
  \item \textbf{下界（式 \ref{eq:bounded-min}）}：当 $\mathit{predictor} > \mathit{GTLI}$ 时，残差最小可以达到 $\mathit{GTLI} - \mathit{predictor}$（一个负数）；但当 $\mathit{predictor} \leq \mathit{GTLI}$ 时，$\mathrm{GCLI} \geq \mathrm{GTLI}$ 意味着残差不能为负，所以下界是 0。
  \item \textbf{上界（式 \ref{eq:bounded-max}）}：$\mathrm{GCLI}$ 最大为 15。如果 $\mathit{predictor} \geq \mathit{GTLI}$，最大残差是 $15 - \mathit{predictor}$；如果 $\mathit{GTLI} > \mathit{predictor}$，$\mathrm{GCLI}$ 的起点是 $\mathit{GTLI}$，最大残差是 $15 - \mathit{GTLI}$。两者合一就是 $15 - \max(\mathit{predictor}, \mathit{GTLI})$。
\end{itemize}

\subsubsection{实现方式}

\begin{lstlisting}[language=C,numbers=none]
// Step 1: 根据 predictor 和 GTLI 计算有效范围
int bounded_code_get_min_max(int predictor, int gtli,
                             int* min_allowed, int* max_allowed)
{
    *min_allowed = -MAX(predictor - gtli, 0);
    *max_allowed = MAX(MAX_GCLI - MAX(predictor, gtli), 0);
}

// Step 2: 将残差映射到紧凑的无符号编码
int bounded_code_get_unary_code(int8_t val, int8_t min_allowed,
                                int8_t max_allowed)
{
    int trigger = abs(min_allowed);
    int aval = abs(val);
    if (aval <= trigger)          // 范围内：正负交替编码
        return (val < 0) ? 2*aval - 1 : 2*aval;
    else                          // 超出 trigger：线性编码
        return trigger + aval;
}
\end{lstlisting}

\subsubsection{Bounded 编码示例}

假设 $\mathit{predictor} = 4, \mathit{GTLI} = 2$：

\begin{lstlisting}[numbers=none]
min_allowed = -MAX(4-2, 0) = -2
max_allowed = MAX(15 - MAX(4, 2), 0) = 11
\end{lstlisting}

有效残差范围 $[-2, 11]$（共 14 个值），比标准 4-clipped 的 25 个值更窄。

对于残差 $\mathit{val} = 0$：$\mathit{aval} = 0 \leq \mathit{trigger} = 2$，编码 $= 2 \times 0 = 0$，unary 编码 \texttt{1}（1 bit）。

对于残差 $\mathit{val} = +2$：$\mathit{aval} = 2 \leq \mathit{trigger}$，编码 $= 2 \times 2 = 4$，unary 编码 \texttt{00001}（5 bit）。

对于残差 $\mathit{val} = +5$：$\mathit{aval} = 5 > \mathit{trigger}$，编码 $= 2 + 5 = 7$，unary 编码 \texttt{00000001}（8 bit）。

\begin{engineeringnote}
同样残差 +5，标准 4-clipped 编码需要 $5 + 4 = 9$ bit，bounded 编码只需 8 bit。当 GTLI 较大时（截断严重），bounded 的优势更明显——因为有效范围更窄，编码值更小。但 bounded 编码需要解码端也已知 predictor 和 GTLI，所以只能在这些信息可用时使用。
\end{engineeringnote}

\subsection{GCLI 方法信号}

每个 band 在 precinct header 中用 3 bit 信号指定 GCLI 编码方法：

\begin{table}[H]
\centering
\caption{GCLI 方法 signaling bits}
\label{tab:gcli-methods}
\begin{tabular}{@{}clll@{}}
\toprule
bits & run & 预测 & 显著性标志 \\
\midrule
0b000 & 无 & 无 & 无 \\
0b001 & 有 & 无 & 无 \\
0b010 & 无 & 垂直 & 无 \\
0b011 & 有 & 垂直 & 无 \\
0b100 & 无 & 无 & 有 \\
0b101 & 有 & 无 & 有 \\
0b110 & 无 & 垂直 & 有 \\
0b111 & 有 & 垂直 & 有 \\
\bottomrule
\end{tabular}
\end{table}

其中 bit 2 = 使用显著性标志，bit 1 = 使用垂直预测。

\section{处理流程}

\subsection{编码端 GCLI 处理}

\begin{lstlisting}[numbers=none]
for each precinct:
    1. update_gclis()
       for each band, each row:
           compute_gcli_buf()  // 按 N_g=4 为一组计算 GCLI

    2. 预测计算
       tco_pred_ver_compute()  // 垂直预测残差

    3. 显著性标志编码
       for each band, each row:
           if method uses sig flags:
               sig_flags_filter_values()
               sig_flags_write()

    4. GCLI 值编码
       for each band, each row:
           if raw fallback:
               pack_raw_gclis()  // 4 bit 直接编码
           else if unary:
               unary_encode()    // unary 编码预测残差
           else if bounded:
               bounded_encode()  // bounded unary 编码
\end{lstlisting}

\subsection{解码端 GCLI 处理}

逆向执行：先解码显著性标志，再解码 GCLI 预测残差，最后通过逆预测恢复 GCLI 值。

\subsection{GCLI 预算计算链路}

编码端在码率搜索时需要\textbf{预算}（不实际编码，只计算编码量）。完整链路：

\begin{lstlisting}[numbers=none]
fill_gcli_budget_table()                    // gcli_budget.c:226
  ├── compute_residuals()                   // gcli_budget.c:168
  │     └── tco_pred_ver() / tco_pred_none() // 对每个 band/row/gtli 计算预测残差
  └── gcli_budget_compute_generic()         // gcli_budget.c:97
        for each position (band row):
          for each gtli (0..MAX_GCLI):
            1. 若使用 sig flags:
               ├── sig_flags_init()
               ├── bgt_sigf = sig_flags_budget()  // 显著性标志的 bit 数
               └── sig_flags_filter_values()       // 过滤掉被跳过的组
            2. GCLI 编码预算:
               ├── 无预测: budget_getunary_unsigned()  // 直接 unary
               ├── 有预测: budget_getunary()           // 带符号 unary
               └── bounded: budget_bounded_code()      // bounded unary
            3. 写入 pbt:
               pbt[position][gtli].sigf = bgt_sigf
               pbt[position][gtli].gcli = bgt_gcli
\end{lstlisting}

\textbf{关键}：\texttt{pbt}（\texttt{precinct\_budget\_table\_t}）是一个二维表，第一维是 position（band 行），第二维是 gtli（量化级别）。每个 cell 记录在该 gtli 下该 position 的 sigf + gcli 编码量。

\subsubsection{源码摘录：\texttt{fill\_gcli\_budget\_table()}}

\begin{lstlisting}[language=C,numbers=none]
int fill_gcli_budget_table(
    uint32_t active_methods,
    const precinct_t* prec,
    const precinct_t* prec_top,
    int* gtli_top_array,
    precinct_budget_table_t* pbt,
    predbuffer_t* residuals,
    int update_only,
    int sigflags_group_width)
{
    // (1) 对每个预测类型/方法/band/row/gtli，计算预测残差
    compute_residuals(active_methods, prec, prec_top,
                      gtli_top_array, residuals);

    // (2) 对每个启用的方法，计算 sigf + gcli 预算
    for (int method = 0; method < GCLI_METHODS_NB; method++)
    {
        if (gcli_method_is_enabled(active_methods, method, ...) == 0)
            continue;
        gcli_budget_compute_generic(prec, method, pbt, residuals,
                                    sigflags_group_width, active_methods);
    }
    return 0;
}
\end{lstlisting}

其中 \texttt{compute\_residuals()}（\texttt{gcli\_budget.c:168}）对每个 band/row/gtli 调用 \texttt{tco\_pred\_ver()} 或 \texttt{tco\_pred\_none()}，计算预测残差并存储在 \texttt{residuals} 缓冲区中。\texttt{gcli\_budget\_compute\_generic()}（\texttt{gcli\_budget.c:97}）遍历所有 position 和 gtli，计算 sigf 和 gcli 预算并写入 PBT。

\subsection{数据预算与 Fs 的影响}

数据（bitplane）预算由 \texttt{fill\_data\_budget\_table()}（\texttt{data\_budget.c:118}）计算，核心函数是 \texttt{budget\_get\_data\_budget()}：

\begin{lstlisting}[language=C,numbers=none]
for each position (band row):
    for gtli = 0..MAX_GCLI:
        n_bitplanes = gcli - gtli
        if group_size == 1:       // 单系数组
            n_bitplanes -= 1      // 符号位不占 bitplane
            if include_sign:      // Fs=1: 独立符号
                n_bitplanes += 1
        if group_size > 1:        // 多系数组 (N_g=4)
            if include_sign:      // Fs=0: 联合符号
                n_bitplanes += 1  // 每个 bitplane 前多一个符号位
        budget[gtli] += n_bitplanes * group_size
\end{lstlisting}

\textbf{Fs=0（联合符号）}：在每个 bitplane 前写入 4 个符号位，因此 \texttt{n\_bitplanes} 多 1。符号位计入数据段。

\textbf{Fs=1（独立符号）}：符号位单独编码（\texttt{budget\_get\_sign\_budget()}），数据段只包含幅度 bitplane。$\texttt{n\_bitplanes} = \texttt{gcli} - \texttt{gtli} - 1$（减 1 是因为 GCLI 本身包含符号位信息）。

符号预算 \texttt{budget\_get\_sign\_budget()}（\texttt{data\_budget.c:90}）逐系数计算：对每个非零系数（逆量化后），累加 1 bit。当系数被量化为零时停止（GTLI 以上的 bitplane 全为零）。

\subsection{流程图}

\begin{figure}[H]
\centering
\begin{tikzpicture}[
    node distance=0.7cm and 0.6cm,
    block/.style={rectangle, draw=structurecolor!60, fill=structurecolor!10,
                  text width=3.6cm, align=center, minimum height=0.7cm,
                  font=\small, inner sep=4pt},
    decision/.style={diamond, draw=second!70, fill=second!12,
                     text width=2.4cm, align=center, aspect=2,
                     font=\small, inner sep=2pt},
    arrow/.style={-{Stealth[length=2.5mm]}, thick, draw=structurecolor!60}
]
    \node[block] (coeff) {系数 sign-magnitude};
    \node[block, below=of coeff] (gcli) {\btt{compute\_gcli\_buf}\\GCLI 值 per 组};
    \node[block, below=of gcli] (pred) {垂直预测残差};
    \node[decision, below=of pred] (sig) {使用 sig\\flags?};
    \node[block, left=2cm of sig] (filter) {\btt{sig\_flags\_filter\_values}\\跳过零残差组};
    \node[block, right=2cm of sig] (all) {全组编码};
    \node[block, below=1.5cm of sig] (encode) {unary/bounded 编码};
    \node[block, below=of encode] (pbt) {GCLI 编码量 $\to$ PBT};

    \draw[arrow] (coeff) -- (gcli);
    \draw[arrow] (gcli) -- (pred);
    \draw[arrow] (pred) -- (sig);
    \draw[arrow] (sig) -- node[above,font=\footnotesize] {是} (filter);
    \draw[arrow] (sig) -- node[above,font=\footnotesize] {否} (all);
    \draw[arrow] (filter) |- (encode);
    \draw[arrow] (all) |- (encode);
    \draw[arrow] (encode) -- (pbt);
\end{tikzpicture}
\caption{GCLI 编码流程图}
\label{fig:gcli-flowchart}
\end{figure}

\section{实现映射}

\begin{implmap}
\begin{tabular}{@{}L{2.2cm}L{3.5cm}L{6.0cm}L{3.8cm}@{}}
\toprule
标准概念 & 入口文件 & 关键函数/结构 & 说明 \\
\midrule
GCLI 计算 & \btt{precinct.c:288} & \btt{compute\_gcli\_buf()} & 按组计算 GCLI \\
GCLI 更新 & \btt{precinct.c:312} & \btt{update\_gclis()} & 遍历所有 band 行 \\
GCLI 预算表 & \btt{gcli\_budget.c:226} & \btt{fill\_gcli\_budget\_table()} & 计算 sigf + gcli 编码量 \\
GCLI 打包 & \btt{packing.c:476} & \btt{pack\_gclis()} & Unary 编码 GCLI \\
GCLI 方法选择 & \btt{precinct\_budget.c:64} & \btt{precinct\_get\_best\_gcli\_method()} & 选最优方法 \\
方法启用 & \btt{gcli\_methods.c:62} & \btt{gcli\_method\_get\_enabled()} & 根据 $R_m$ 选择 run 模式 \\
方法信号 & \btt{gcli\_methods.c:119} & \btt{gcli\_method\_get\_signaling()} & 3 bit 编码 \\
显著性标志 & \btt{sig\_flags.c} & \btt{sig\_flags\_write()}, \btt{sig\_flags\_read()} & 显著性编码/解码 \\
Unary 编码 & \btt{packing.c:271} & \btt{unary\_encode()} & 带符号 unary \\
预测 & \btt{pred.c} & \btt{tco\_pred\_ver\_compute()}, \btt{tco\_pred\_ver\_inverse()} & 垂直预测 \\
\bottomrule
\end{tabular}
\end{implmap}

\begin{codefact}
本章关键结论依据以下函数：
\begin{itemize}
  \item \texttt{compute\_gcli\_buf()} --- \texttt{precinct.c:288} --- GCLI = BSR(or\_all) + 1
  \item \texttt{update\_gclis()} --- \texttt{precinct.c:312} --- 遍历 band 行更新 GCLI
  \item \texttt{fill\_gcli\_budget\_table()} --- \texttt{gcli\_budget.c:226} --- GCLI 预算计算总入口
  \item \texttt{gcli\_budget\_compute\_generic()} --- \texttt{gcli\_budget.c:97} --- sig flags 过滤 + 预算累加
  \item \texttt{budget\_getunary()} --- \texttt{budget.c:102} --- unary 编码预算
\end{itemize}
\end{codefact}

\section{算例}

详见第 \ref{ex:gcli-gtli-budget} 章（GCLI、GTLI 与预算算例），包含：
\begin{itemize}
  \item 例 1：极小例子——从 4 个系数到 GCLI
  \item 例 2：GCLI 编码全流程（预测 $\to$ sigflag $\to$ unary）
  \item 例 2b：Unary 编码详细示例
  \item 例 2c：Bounded 编码示例
  \item 例 3：GTLI 表计算
  \item 例 4：预算表计算
  \item 例 5：GTLI 对编码量的影响
\end{itemize}

\section{常见误解}

\begin{misunderstanding}
\begin{enumerate}
  \item \textbf{``GCLI 是每个系数的位深''} --- 不是。GCLI 是\textbf{每组 4 个系数}共享的，表示该组的最大位深。
  \item \textbf{``GCLI 和 GTLI 是同一回事''} --- 不是。GCLI 从系数计算得出（描述``这组有多大''），GTLI 由码率控制决定（描述``预算只够保留到哪一位''）。$\mathrm{GCLI} \geq \mathrm{GTLI}$，差值越大说明量化越狠。
  \item \textbf{``Unary 编码总是比 raw 编码效率高''} --- 不一定。当 GCLI 值较大且分布均匀时，raw 4-bit 编码可能更高效。Raw fallback 机制正是为此设计。
  \item \textbf{``显著性标志只用于全零组''} --- ZRCSF 模式确实只跳过全零组；但 ZRF 模式跳过的是预测残差为零的组（即使系数不为零）。
  \item \textbf{``预测的是系数值''} --- 不是。预测的是 \textbf{GCLI 值}（每组的位深摘要），不是单个系数。残差 = 当前 GCLI $-$ 上方 GCLI。
\end{enumerate}
\end{misunderstanding}

\section{与前后章衔接}

\begin{itemize}
  \item \textbf{输入}：本章的输入来自第 \ref{ch:bands-precincts} 章（Precinct）提取的 sign-magnitude 系数和 GCLI 值
  \item \textbf{输出}：本章的输出（GCLI 编码量预算）写入 \texttt{precinct\_budget\_table\_t}，供第 \ref{ch:rate-control} 章（码率控制）的 Q/R 搜索使用
\end{itemize}

\section{本章小结}

GCLI 是 JPEG XS 轻量熵编码的核心。它回答一个简单的问题：``这组系数最多用到了第几个 bitplane？'' 通过 OR 运算计算每组系数的位深摘要，然后利用相邻 precinct 的相关性（垂直预测），只编码一个小的残差值。结合显著性标志跳过零组/零残差组，再用 unary 编码压缩残差，JPEG XS 实现了无需算术编码的高效系数位深编码。

初学者需要记住的核心链条：\textbf{系数 $\to$ OR $\to$ BSR+1 = GCLI $\to$ 减去预测值 $\to$ 残差 $\to$ sigflag 过滤 $\to$ unary 编码}。

下一章将讲解预算表和码率控制——如何根据总比特预算决定每个 precinct 的量化强度（GTLI）。
